% called by main.tex
\chapter{Sprint 1 : Image Uploading , Processing And Scanning }
\label{ch::chapter4}

\section{Introduction}
In this chapter, we will discuss the first Sprint of our project in details. We will begin by presenting the sprint backlog. Next, we will dive into the design phase, presenting the necessary diagrams needed to illustrate our approach. Finally, we will cover the development phase, detailing how each feature was implemented.

\section{Sprint backlog}
The table 3.1 below presents the sprint 1 backlog.
\begin{table}[!htpb]
    \caption{Sprint 1 backlog.}
    \centering
    \begin{tabular}{|>{\centering\arraybackslash}m{0.05\linewidth}|>{\centering\arraybackslash}m{0.15\linewidth}|>{\centering\arraybackslash}m{0.45\linewidth}|>{\centering\arraybackslash}m{0.1\linewidth}|>{\centering\arraybackslash}m{0.15\linewidth}|} \hline 
    \textbf{ID} & \textbf{Functionality} & \textbf{User Story} & \textbf{Priority} & \textbf{Complexity} \\ \hline
         1 & Image Upload & As a user, I want to be able to upload one or more images of different types like JPEG or PNG so that I can process various documents easily. & High & Medium\\ \hline 
         2 & Image Preprocessing & As a developer, I want the system to enhance the quality of uploaded images by resizing, reducing noise, adjusting contrast, grayscale conversion and so that the user receives more accurate results. & High & High\\ \hline
         3 & OCR Integration & As a user, I want the system to read text from images using OCR so that I can extract information without typing it manually. & High & High\\ \hline
         4 & HMI & As a user, I want the interface to be user-friendly and intuitive, allowing me to upload images, view and edit text, and export data without any confusion or difficulty. & Medium & Low\\ \hline
    \end{tabular}
    
    \label{tab:sprint1_backlog}
\end{table}
\newpage


\section{Sprint functionalities}
In this section, we will discuss each of the current sprint's specific functionalities .

\subsection{Image upload}
In this phase of the project, the user will be able to either upload a single image or multiple images of different document types.

\subsection{Image preprocessing}
The application moves the submitted images through the Image Preprocessing phase, which is series of operations using OpenCV designed to optimize them for accurate and efficient OCR scanning.

\subsection{OCR scan}
The images then get to the OCR stage, where we use Tesseract OCR to extract text from them.

\subsection{HMI}
The user-friendly HMI allows the user to use all the application's functionalities without any difficulty or confusion through clear visual elements and straightforward error messages.\\


\section{Design}
This section is where we will presents the architecture design of the current sprint, which includes the static and dynamic diagrams.\\
\subsection{Static diagrams}
In this section we will display the static architecture design of the sprint 1.
\subsubsection{Use case diagram}
The figure 3.1 presents the use case diagram of the sprint 1.
\vspace{2mm}
\begin{figure}[!htpb]
    \centering
    \includegraphics[width=10cm]{figures/Sprint 1 Use Case.png}
    \caption{Sprint1 Use Case Diagram.}
    \label{fig:Sprint1_use_case}
\end{figure}

\subsubsection{Class diagram}
The figure 3.2 presents the class diagram of the sprint 1.
\begin{figure}[!htpb]
    \centering
    \includegraphics[height=5cm]{figures/Sprint 1 Class.png}
    \caption{Sprint1 class diagram.}
    \label{fig:Sprint1_Class}
\end{figure}

\subsection{Dynamic diagrams}
In this section we will display the dynamic architecture design of the sprint 1
\subsubsection{Sequence diagram}
The sequence diagram of this sprint shows the changes that occur within our system over time between the user and the two other entities, the user interface and the application from the image upload to the OCR phase, as shown in the figure 3.3 below.
\begin{figure}[!htpb]
    \centering
    \includegraphics[height=10.5cm]{figures/Sprint 1 Sequence.drawio.png}
    \caption{Sprint1 sequence diagram.}
    \label{fig:Sprint1_sequence}
\end{figure}
\newpage
\section{Development}
In this section, we detail the realization of each sprint functionalities, sharing the technical steps behind our OCR application, in addition to showcasing the output and UI of each functionality through some screenshots. Starting with the first sprint.
\subsection{Image upload}
Through the image upload functionality, the user can select to upload either one image or multiple images then click the submit button, it first goes through the client-side verification, which is the JavaScript code that checks how the existence of selected images and if they have the appropriate extensions.\\
It is possible for the user to upload any of the most common image extensions, which are : \\
.xbm, .tif, .jfif, .ico, .tiff, .gif, .svg, .jpeg, .svgz, .jpg, .webp, .png, .bmp, .pjp, .apng, .pjeg and .avif.\\
Once The user selected the images, he can press on the Submit button.\\
The application then receives a POST request containing the image files, the images then get processed sequentially on the server-side, where the server checks for the file format and size to ensure it meets the requirements. Finally, the image gets saved to the images folder in media root, which is the path defined in settings.py file.

The figure 3.4 shows the uploading interface.

\vspace{10mm}
\begin{figure}[!htpb]
    \centering
    \includegraphics[width=15cm]{figures/Image_Upload_Interface.png}
    \caption{Image upload interface}
    \label{fig:Image_Upload_Interface}
\end{figure}
The figure 3.5 displays the error message when the user clicks the submit button without importing any image.
\begin{figure}[!htpb]
    \centering
    \includegraphics[width=15cm]{figures/Image_Upload_Interface_error1(highlighted).png}
    \caption{Image upload interface error}
    \label{fig:Image_Upload_Interface_error1}
\end{figure}

The figure 3.6 displays the error message when the user try to import image forms that are not supported by our system.
\begin{figure}[!htpb]
    \centering
    \includegraphics[width=15cm]{figures/Image_Upload_Interface_error2(highlighted).png}
    \caption{Image upload interface error2}
    \label{fig:Image_Upload_Interface_error2}
\end{figure}

\newpage

\subsection{Image preprocessing}
The application then performs the same preprocessing algorithm on all the uploaded images in iterations, where it loads the image from the media root and reads some of its properties like height, width and path.
The image is then converted to grayscale to make the text more highlighted for better OCR detection, then reducing Noise for more accuracte results. Finally, estimating the text size to dynamically adjust parameters for adaptive thresholding, which uses them to achieve the best foreground-to-background ratio for proper image binarization. The preprocessed version now gets saved directly to the media root under the name "preprocessed\textunderscore id".\\
The figures 3.7 and 3.8 displays the two versions of the same image, before and after the Preprocessing phase.

\vspace{10mm}
\begin{figure}[!htpb]
    \centering
    \begin{minipage}{0.48\textwidth}
        \centering
        \includegraphics[height=11.5cm]{figures/Preprocessing1.jpg}
        \caption{Uploaded Image.}
        \label{fig:Preprocessing1}
    \end{minipage}
    \hfill
    \begin{minipage}{0.48\textwidth}
        \centering
        \includegraphics[height=11.5cm]{figures/Preprocessing2.jpg}
        \caption{Preprocessed image.}
        \label{fig:Preprocessing2}
    \end{minipage}
\end{figure}

\subsection{OCR scan}
The images then go through the OCR function, where we use Tesseract OCR to extract text from each Preprocessed image and clean the text from unnecessary spaces, linebreaks and characters which aren't commonly used in documents in order to make the text less noisy.

As an example here's how the extracted text of the figure 3.7 above looks as shown in the figure 3.9 .

\begin{figure}[!htpb]
    \centering
    \includegraphics[width=17cm]{figures/Tesseract_output.png}
    \caption{OCR initial output.}
    \label{fig:Tesseract_output}
\end{figure}
\subsection{HMI Development}
The user-friendly HMI allows the user to use all the application's functionalities without any difficulty or confusion through clear visual elements and straightforward error messages.\\
On the first interface, the user will find the project's logo and description at the top, followed by a guide section providing simple instructions on how to use it. Below this, there is a section containing two buttons for selecting either a single image or multiple images and a third one for submitting the selected image(s). The Submit button also redirects the user to the second interface after a brief moment of waiting until the Preprocessing, OCR, Document type detection, and text organization processes are completed.\\
The second interface contains a modernly designed container in which the user will find the names of the selected images, their document types, and extracted data clearly organized. Below this container, the user will find two buttons: one for editing the extracted data to ensure it is error-free and the other for submitting the changes and displaying the export functionality container. Through the export functionality container, the user can select at least one format from the checkboxes labeled as CSV, JSON, and WORD, then press the export button to download the data onto their machine in the selected formats.




\section{Conclusion}
In conclusion, Sprint 1 marks the initial stage of constructing OCR app, which contains some of its key features such as Image upload, preprocessing and text extraction. This sprint is essential as it forms a strong foundation for the project's future development.
